\section{NsJail e Container}\label{nsjail-e-container}

Si vuole in questo capitolo approfondire i meccanismi utilizzati da
NsJail per l'isolamento delle risorse del kernel Linux.\footnote{\label{fn:ref_30}Google
  Security Team, \emph{NsJail Process Isolation Tool},
  https://nsjail.dev/} Tutti questi meccanismi utilizzati da NsJail sono
utilizzati anche da Docker e Incus per implementare i container e verrà
visto brevemente in che modo.

\subsection{Namespaces}\label{namespaces}

Un namespace è una feature del kernel Linux che permette di partizionare
le risorse del kernel in modo che:\footnote{\label{fn:ref_31}Linux manual page,
  \emph{namespaces(7) - Overview of Linux namespaces},
  https://man7.org/linux/man-pages/man7/namespaces.7.html}

\begin{itemize}
\tightlist
\item
  un insieme di processi è in grado di vedere un certo insieme di
  risorse
\item
  un altro insieme di processi è in grado di vedere un altro insieme di
  risorse.
\end{itemize}

Il kernel Linux supporta diversi tipi di namespace: Cgroup, IPC,
Network, Mount, PID, Time, User ed UTS.

Il \textbf{PID namespace} permette di partizionare gli ID dei processi
in modo che per esempio, su due PID namespace differenti possano
coesistere due processi con process ID identici. Inoltre un processo che
esegue sul pid namespace \passthrough{\lstinline!A!} non può vedere
quali processi sono in esecuzione sul pid namespace
\passthrough{\lstinline!B!}.

Il \textbf{Network namespace} invece offre l'isolamento delle risorse di
networking, come device, lo stack IPv4, le tabelle di routing, ecc. Un
device fisico vive all'interno di un unico network namespace, e le
connessioni tra namespace sono gestite dai device
\passthrough{\lstinline!veth!}.\footnote{\label{fn:ref_32}Linux manual page,
  \emph{network\_namespaces(7) - Network namespaces},
  https://man7.org/linux/man-pages/man7/network\_namespaces.7.html}

In sintesi, i container sfruttano i namespace per creare un ambiente
isolato in cui eseguire i processi, montare il filesystem e configurare
le interfacce di rete.

\subsection{Seccomp}\label{seccomp}

Seccomp permette a un processo di specificare un filtro per le chiamate
di sistema.\footnote{\label{fn:ref_33}Linux manual page, \emph{seccomp(2) - Operate on
  Secure Computing state of the process},
  https://man7.org/linux/man-pages/man2/seccomp.2.html} Il filtro è
espresso come un programma \passthrough{\lstinline!cBPF!} (classic BPF),
che opera sul numero della syscall e sui suoi argomenti. Fornisce un
meccanismo essenziale per minimizzare la superficie di attacco verso il
kernel.

A titolo esemplificativo, l'utilizzo dell'API
\passthrough{\lstinline!libseccomp!} in linguaggio
\passthrough{\lstinline!C!} prevede la seguente sequenza operativa:

\begin{itemize}
\tightlist
\item
  Si inizializza nel codice il contesto del filtro:
  \passthrough{\lstinline!scmp\_filter\_ctx ctx!}
\item
  Si aggiungono le regole selettive invocando
  \passthrough{\lstinline!seccomp\_rule\_add(ctx, ...)!}
\item
  Si caricano le regole nel kernel con
  \passthrough{\lstinline!seccomp\_load(ctx, ...)!}. Da questo istante,
  il processo e i suoi discendenti saranno soggetti ai vincoli definiti.
\end{itemize}

I container Docker tradizionali applicano Seccomp per bloccare le
chiamate di sistema considerate pericolose, mentre Incus nei system
container privilegiati o semi-privilegiati adotta una policy permissiva
delegando l'isolamento agli user namespace.

\subsection{cgroups}\label{cgroups}

I Control Groups costituiscono un meccanismo per aggregare o
partizionare insiemi di task e i loro discendenti in gerarchie
controllate.\footnote{\label{fn:ref_34}The Linux Kernel Organization, \emph{Control
  Groups (cgroups) Documentation},
  https://www.kernel.org/doc/html/latest/admin-guide/cgroup-v1/cgroups.html}

In particolare, la documentazione del kernel Linux formalizza i seguenti
concetti:

\begin{itemize}
\tightlist
\item
  \passthrough{\lstinline!cgroup!}: associa un insieme di task a un
  insieme di parametri e limiti per uno o più sottosistemi.
\item
  \passthrough{\lstinline!sottosistema!}: noto anche come \emph{resource
  controller}, schedula una risorsa hardware o applica quote e vincoli a
  un determinato raggruppamento di task.
\end{itemize}

Tramite un \passthrough{\lstinline!cgroup!} è possibile ad esempio
limitare la quota di CPU o di memoria RAM assegnata a un gruppo di task.
Nei container, i \passthrough{\lstinline!cgroup!} sono utilizzati
sistematicamente per imporre tetti massimi all'utilizzo di risorse
computazionali e condurre attività di accounting e monitoraggio.
